\documentclass[10pt,a4paper]{amsart}
\usepackage[margin=19mm]{geometry}
\usepackage{amsmath,amssymb,mathtools}
\usepackage[scr=rsfs,cal=euler]{mathalfa}
\usepackage{xfrac}
\usepackage[unicode,hidelinks,bookmarks=false]{hyperref}
\newcommand{\ie}{i.e.\ }
\newcommand{\cf}{cf.\ }
\newcommand{\resp}{respectively\ }
\newcommand{\Subsection}{\subsection}
\providecommand{\nobreakdash}{\nobreak\hbox{-}}
\setlength{\emergencystretch}{2em}
\multlinegap0pt
\usepackage{amssymb}
\usepackage[shortlabels]{enumitem}
\usepackage{dsfont}
\usepackage{xcolor}
\newtheorem{proposition}{Proposition}
\newtheorem{lemma}{Lemma}
\newtheorem{remark}{Remark}
\def\G{\Gamma}
\newcommand{\N}{\numberset{N}}
\newcommand{\RCA}{\mathsf{RCA}_0}
\def\WO{{\textrm{WO}}}
\def\a{\alpha}
\newcommand{\andrea}[1]{\textcolor{green}{#1}}
\def\b{\beta}
\def\conc{{}^\smallfrown}
\def\d{\delta}
\def\g{\gamma}
\newcommand{\gd}[1]{{\ulcorner{#1}\urcorner}}
\newcommand{\homo}{homogeneous}
\def\ifff{\Longleftrightarrow}
\def\la{\langle}
\newcommand{\leo}{<\sb{o}}
\newcommand{\leqo}{\leq\sb{o}}
\newcommand{\lorenzo}[1]{\textcolor{blue}{#1}}
\newcommand{\numberset}{\mathbb}
\newcommand{\pluso}{+\sb{o}}
\def\ra{\rangle}
\newcommand{\set}[1]{{\{{#1}\}}}
\def\stops{{\downarrow}}
\title{The strength of Ramsey's theorem for $\a$-large sets}
\author{Lorenzo Carlucci, Andrea Volpi, Konrad Zdanowski}
\date{Source: arXiv:2603.22579v2 (CC BY 4.0)}
\begin{document}
\maketitle

\noindent\emph{Source and licence.} The strength of Ramsey's theorem for $\a$-large sets. Version arXiv:2603.22579v2 (CC BY 4.0), distributed by arXiv under the Creative Commons Attribution 4.0 International licence, http://creativecommons.org/licenses/by/4.0/, as recorded in the arXiv licence metadata for this exact version and shown on the abstract page https://arxiv.org/abs/2603.22579v2. Full licence notice: \texttt{licenses/arxiv-2603.22579-CC-BY-4.0.txt}.\par\smallskip

\noindent\textit{Editorial scope.} Counted unit: the article's lemma lem:lem:eq:homo:wit (source0.tex lines 1189--1197). The article's complete original proof of it is delivered with it (source0.tex lines 1199--1310, one \texttt{proof} environment of 112 lines). No mathematics is written, repaired or re-derived: the statement and the proof are the article's own lines, in the article's own notation, and no retained line is substituted or re-flowed.  Retained uncounted as the source-local context the proof needs: lines 523--527; lines 736--742; lines 816--820; lines 872--880; lines 1098--1103.  Nothing was omitted from any retained span, so the package carries no omission record.  No retained line contains a citation, so the wrapper carries no bibliography and no bibliography entry.  No retained line contains a figure environment, an image-inclusion command or drawing code, so no image is delivered and the package declares no asset dependency.  Editorial formatting only, in addition: an AMS \texttt{amsart} preamble that restates the article's own declarations and macros used by the retained lines, namely the environments \texttt{proposition}, \texttt{lemma}, \texttt{remark} on separate counters, together with the standard packages those lines need.  The retained environments print in this document as Proposition 1, Lemma 1, Remark 1 in order of appearance, whereas the article prints the same items with the numbering of its own full document.  The complete native source of the article as distributed in the arXiv e-print for this exact version is bundled unchanged as \texttt{sources/197029/source0.tex}.
\begin{proposition}\label{prop: goodnorm}
For all $\b, \d < \G_0$, if $\b < \d$ then $\b \le \d[|\b|]$.
Moreover, if $s \in [\N]^{<\omega}$, $\d[s] \le \b$ and $\min s \ge |\b|$ then for some $t \sqsubseteq s$ we have $\d[t] = \b$.
Consequently, if $s$ is $\d$-large and $\min s \ge |\b|$ then $s$ is $\b$-large.
\end{proposition}

\begin{lemma}[$\RCA+\WO(\Lambda)$]\label{lem:fseq:bounding}
Let $\alpha\in\Lambda$, let $x\sb0,\dots,x\sb n\in \N$ and let $\gd{\alpha}<x\sb0 <\dots< x\sb n$.
For $i=0,\dots,n+1$, let $\alpha\sb i= \a[x\sb 0, \dots, x\sb{i-1}]$.
Then, for each $i<n$, $\gd{\alpha\sb i}< x\sb i$.
\andrea{comment: with the previous notation the brackets denoted the code of the ordinal which was equal to its pseudonorm. now they are not equal anymore and I don't understand if we care about the pseudonorm or the code. I suspect it is the code}
\lorenzo{seems to be the code}
\end{lemma}

\begin{remark}\label{pairingfunctionproperty}
We fix a computable pairing function $\langle \cdot, \cdot\rangle : \N\times \N\to \N$. 
Notice that we can assume without loss of generality that $|\gamma| \le \la \gd\g,z \ra$ for all $\gamma \in \Lambda$ and $z\in\N$. 
We write $\langle \g, z\rangle$ for $\langle \gd\g, z\rangle$.
\end{remark}

\begin{lemma}[$\RCA+\WO(\Lambda)$]\label{lem:from:alpha:to:beta}
Let $\beta, \alpha\in\Lambda$ with $\beta\leqo \alpha$.
Let $t$ be a $\beta$-size sequence and let $s$ be such that $s < t$ and $\alpha[s]=\beta$.
Then 
\begin{multline*}
\set{\langle\gamma,z\rangle\colon M\sb\beta(\langle\gamma,z\rangle,t)\mbox{ accepts and }\gamma\leqo \beta } =\\
\set{\langle\gamma,z\rangle\colon M\sb\alpha(\langle\gamma,z\rangle, s \conc t)\mbox{ accepts and }\gamma\leqo \beta }.
\end{multline*}
\end{lemma}

\begin{lemma}[$\RCA+\WO(\Lambda)$]\label{lem:bounding:computations}
Let $\alpha\in\Lambda$ and let $H\subseteq\N$ be an infinite $c\sb{\alpha \pluso3}$-\homo\ set, let $\beta\leqo \alpha$, let $s$ be a $\beta$-size sequence in $S\sp2(\alpha,H)\sp{>n_{\a,\b}}$ and let $s\sb0, s\sb1$ be in $S\sp2(\alpha,H)\sp{>n_{\a,\b}}$ such that $s\sb0 < s\sb1 < \min s$. 
Then, 
$$\forall e,x<s\sb0 \left(\set{e}\sp {Y \sb \beta \sp s}(x)\stops\sb {\min s}\Longrightarrow \set{e}\sp {Y \sb \beta \sp s}(x)\stops\sb {s\sb1}\right),$$
where $Y_\beta^s =\set{y<\max s\colon M\sb\beta(y,s)\mbox{ accepts}}$. 
\end{lemma}

\begin{lemma}[$\RCA+\WO(\Lambda)$]\label{lem:lem:eq:homo:wit}
Let $H\subseteq\N$ be an infinite $c\sb{\alpha \pluso3}$-\homo\ set. 
Then, for each $\beta\leqo \alpha$ and for each choice of $\beta$-size sequences $s,t$ in $S\sp2(\alpha,H)\sp{>n_{\a,\b}}$ such that $s < t$, the following equivalence holds
$$
M\sb\beta(\langle \gamma,e\rangle,s)\mbox{ accepts} \ifff
M\sb\beta(\langle \gamma,e\rangle,t)\mbox{ accepts},
$$
for all $\langle \gamma,e\rangle<\min s$.
\end{lemma}

\begin{proof}
Let us assume that we have $\beta\leqo\alpha$
and $k, n \in \N^+$ such that $s=\langle a\sb0,\dots,a\sb {k-1}\rangle$ and  $t=\langle b\sb0,\dots, b\sb {n-1}\rangle$ from $S\sp2(\alpha,H)\sp{>n_{\a,\b}}$ satisfying the hypotheses of the Lemma and such that they falsify the thesis.
First, let us observe that for all $\gamma\leqo \beta$ such that $|\gamma | < \min s$, there are $k'\leq k$ and $n'\leq n$ such that 
$$
\gamma=\beta[a\sb0]\cdots [a\sb{k'-1}]=\beta[b\sb0]\cdots[b\sb{n'-1}].
$$
Indeed, if $|\gamma| < \min s < \min t$ (by Remark \ref{pairingfunctionproperty} about the choice of the pairing function) and $\gamma\leqo \beta$ then, by Proposition \ref{prop: goodnorm}, $\gamma$ has to occur as $\beta[a\sb0]\cdots [a\sb{k'-1}]$ and as $\beta[b\sb0]\cdots [b\sb{n'-1}]$ for some $k'\leq k$ and $n'\leq n$. 

\iffalse
\lorenzo{what is the need for $\delta$ below here?? I moved this a few lines down where it seems needed}
Moreover, if $\gamma=\beta[a\sb0,\dots, a\sb{k'-1}]$ and $|\delta|\leq a\sb{k'}$ is such that $\delta\leqo\gamma$ then there exists $k''$ such that $k'\leq k''\leq k$ and
$$
\delta=\beta[a\sb0,\dots, a\sb{k''-1}].
$$
Indeed, since $\delta\leq a\sb{k'}$  we may again use Proposition \ref{prop: goodnorm}. 
\fi

So, let $\gamma\leqo\beta$ and $k'\leq k$ and $n' \leq n$ be such that 
$$
\gamma=\beta[a\sb0]\cdots [a\sb{k'-1}]=\beta[b\sb0]\cdots [b\sb{n'-1}].
$$
The numbers $k'$ and $n'$ are unique with respect to $\gamma$ and to the sequences $s$ and $t$ so we will denote them by $k\sb\gamma$ and $n\sb\gamma$ respectively. 
We also set $k\sb\beta=n\sb\beta=0$ as $\beta[\emptyset]=\beta$. 

%We will need that, by Lemma \ref{lem:fseq:bounding}, for each $\gamma$ with $k\sb\gamma$ properly defined, we have that $\gamma<a\sb{k\sb\gamma+1}$. \lorenzo{is actually not needed below!!}

Let us observe that for all $\gamma\leqo\beta$ such that $k\sb\gamma$ and $n\sb\gamma$ are defined and for all $\la\delta,e\ra<a\sb{k\sb\gamma}$, again by Lemma \ref{lem:from:alpha:to:beta} we have the following equivalences.
$$
M\sb\beta(\la\delta,e\ra,s)\mbox{ accepts} \ifff          M\sb\gamma(\la\delta,e\ra,\langle a\sb{k\sb\gamma},\dots, a\sb k \rangle)\mbox{ accepts}
$$
and
$$
M\sb\beta(\la\delta,e\ra,t)\mbox{ accepts}\ifff  M\sb\gamma(\la\delta,e\ra,\langle b\sb{n\sb\gamma},\dots, b\sb n\rangle)\mbox{ accepts}.
$$
For all $\delta \leqo\gamma$ such that $|\delta|< a\sb{k\sb\gamma}$ (by Remark \ref{pairingfunctionproperty} about the choice of the pairing function) the numbers $k\sb\delta$ and $n\sb\delta$ are also well defined. 
Indeed, if $\gamma=\beta[a\sb0]\cdots [a\sb{k'-1}]$ and $|\delta|\leq a\sb{k'}$ is such that $\delta\leqo\gamma$ then there exists $k''$ such that $k'\leq k''\leq k$ and $\delta=\beta[a\sb0]\cdots [a\sb{k''-1}]$. 
Indeed, since $|\delta|\leq a\sb{k'}$  we may again use Proposition \ref{prop: goodnorm}.

By the above observations we thus rule out the case in which $M\sb\beta(\la \delta,e \ra,t)$ accepts some $\la\delta,e\ra$ below $a\sb{k\sb\gamma}$ while $\delta$ does not occur as $\beta [a\sb0]\cdots [a\sb{k'-1}]$ for some $k'\leq k$.

For each $\gamma$ such that $\gamma \leqo\beta$ and such that $k\sb\gamma=k\sb\gamma(s)$ and $n\sb\gamma=n\sb\gamma(t)$ are defined let us define sets
\begin{eqnarray*}
Y\sp {\gamma}\sb{s}&=&
\set{\la\delta,e\ra<a\sb{k\sb\gamma}\colon M\sb\gamma(\la\delta,e\ra,\langle a\sb{k\sb\gamma},\dots,a\sb k\rangle)\mbox{ accepts}},\\
Y\sp{\gamma}\sb{t}&=&      \set{\la\delta,e\ra<b\sb{n\sb\gamma }\colon M\sb\gamma(\la\delta,e\ra,\langle b\sb{n\sb\gamma},\dots,b\sb n\rangle)\mbox{ accepts}}.
\end{eqnarray*}
The above sets are uniformly $\Delta^0_1$-definable with a parameter bounding computations of machines $M\sb\gamma (y,\langle b\sb{n\sb\gamma},\dots,b\sb n\rangle)$ on finite sets of inputs.

We claim that for each $\gamma \leqo \beta$ such that $k_\gamma \leq k$ the sets $Y\sp {\gamma}\sb{s}$ and $Y\sp{\gamma}\sb{t}$ coincide on all $\langle \delta, e\rangle < a_{k_\gamma}$.
Towards a contradiction, we define the set of indices $D$ on which the corresponding $Y\sp {\gamma}\sb{s}$ and $Y\sp{\gamma}\sb{t}$ differ:
\begin{multline*}
%\begin{split}
D=\{k'\leq k\colon \exists \gamma\leq a\sb{k'} 
(k'=k\sb\gamma  \land  \\
\exists \delta<a\sb{k\sb\gamma}\exists e<a\sb{k\sb\gamma}(\delta\leqo\gamma\land
\la\delta,e\ra<a\sb {k\sb\gamma} \land  \\
\neg(\la\delta,e\ra\in Y\sp{\gamma}\sb{s} \ifff \la\delta,e\ra\in Y\sp{\gamma}\sb{t}))\}.
%\end{split}
\end{multline*}
The set $D$ is $\Delta^0_1$-definable with a sufficiently large parameter. Suppose by way of contradiction that $D$ is nonempty.
%\andrea{I agree that $D$ is nonempty because of the assumption of contradiction of the lemma, but is it really true $0 \in D$?}. \lorenzo{I agree just stating that $D$ is non-empty}

Let $k\sb0=k\sb{\gamma\sb0}$ be a maximal element of $D$ and let $n\sb0=n\sb{\gamma\sb0}$.
By the construction of the machine $M\sb{1}$,  $1\leo \gamma\sb0$ and $\gamma\sb0$ has to be a successor ordinal.
Otherwise, if $\gamma\sb0$ is limit then we would also have $k\sb0+1\in D$ contradicting the maximality of $k\sb0$. 
Thus, let $\gamma\sb0=\gamma\sb 1\pluso 1$ and let $\la\delta,e\ra<a\sb{k\sb0}$ be witnesses for $k\sb0\in D$, that is $\la\delta,e\ra<a\sb{k\sb0}$ and 
$$
\neg(\la\delta,e\ra\in Y\sp{\gamma\sb0}\sb{s} \ifff \la\delta,e\ra\in Y\sp{\gamma\sb0}\sb{t}),
$$
where $Y\sp{\gamma\sb 0}\sb{s}$ and $Y\sp{\gamma\sb 1}\sb{t}$ are defined as above. 

We claim that $\delta=\gamma\sb0$. 
Otherwise, we would have $k\sb{\delta}\in D$ for $k\sb\delta>k\sb{\gamma\sb0}$, as $k\sb\delta$ is well defined. Note that $\delta = \gamma_0$ implies $\langle \gamma_0, e\rangle < a_{k_0}$.

We distinguish two cases and derive a contradiction in both of them. 



\smallskip

\textbf{Case 1:} $\langle\gamma\sb0,e\rangle\in Y\sp{\gamma\sb0}\sb{t}$. 
Then
$$
\set{e}\sp{Y\sp{\gamma\sb 1}\sb{t}}(0)\stops\sb{b\sb{n\sb0+1}}.
$$
We need to prove that 
$$
\set{e}\sp{Y\sp{\gamma\sb 1}\sb{s}}(0)\stops\sb{a\sb{k\sb0+1}}.
$$
Note that $\gamma\sb 1 < \gamma\sb 1\pluso 1$ by properties of the system of fundamental sequences and $\gamma\sb 1\pluso 1 = \gamma_0 <a\sb{k\sb0}$ as observed above. Moreover, $\langle b\sb{n\sb0+1},\dots, b\sb n\rangle$ is a $\gamma_1$-size sequence, and $a\sb{k\sb0} < a\sb{k\sb0+1} < b\sb{n\sb0+1}$ since $s<t$ by hypotheses of the Lemma. 
Therefore we can apply Lemma \ref{lem:bounding:computations} and conclude that
$$
\set{e}\sp {Y\sp{\gamma\sb 1}\sb{t}}(0)\stops\sb{a\sb{k\sb0+1}}.
$$
We argue that each query of type $\langle\delta,z\rangle$ asked during the computation of $\set{e}\sp {Y\sp{\gamma\sb 1}\sb{t}}(0)\stops\sb{a\sb{k\sb0+1}}$ is answered in the same way by the oracles $Y\sp{\gamma\sb 1}\sb{t}$ and $Y\sp{\gamma\sb 1}\sb{s}$. 
Indeed, let us consider $\langle\delta, z\rangle<a\sb{k\sb0+1}$. 
We may assume that $\delta\leqo\gamma\sb 1$ since otherwise $\langle\delta,z\rangle$ belongs neither to $Y\sp{\gamma\sb 1}\sb{t}$ nor to $Y\sp{\gamma\sb 1}\sb{s}$. 
Then $\delta<a\sb{k\sb0+1}$ and there are $k\sb\delta$ and $n\sb\delta$ such that 
$$
\delta=\gamma\sb 1[a\sb{k\sb0+1}]\cdots [a\sb{k\sb\delta-1}]=
\gamma\sb 1[b\sb{n\sb0+1}]\cdots [b\sb{n\sb\delta-1}].
$$
By the maximality of $k\sb0$ we get that $\langle\delta,z\rangle\in Y\sp{\gamma\sb 1}\sb{t}$ if and only if $\langle\delta, z\rangle\in Y\sp{\gamma\sb 1}\sb{s}$.
So, we proved that $\la \gamma\sb0,e\ra\in Y\sp{\gamma\sb0}\sb{s}$, contradicting the choice of $\langle\gamma\sb0,e\rangle$.
\smallskip

\textbf{Case 2:} $\langle\gamma\sb0,e\rangle\in Y\sp{\gamma\sb0}\sb{s}$, so
$\set{e}\sp {Y\sp{\gamma\sb 1}\sb{s}}(0)\stops\sb{a\sb{k\sb0+1}}$.
Then, we argue as in the previous case that 
$\set{e}\sp{ Y\sp{\gamma\sb 1}\sb{t}}(0)\stops\sb{a\sb{k\sb0+1}}$ and, therefore, $\la\gamma\sb0,e\ra \in Y\sp{\gamma\sb0}\sb{t}$, contradicting again the choice of~$\langle\gamma\sb 0,e\rangle$.
\end{proof}

\end{document}
